\subsection{The harmonic step only needs span below one}\label{v9:harmonic:section}
For a real number $x$, write $\{x\}=x-\lfloor x\rfloor$.

\begin{lemma}[Extended conditional-span domain]\label{v9:harmonic:span}
Let $Z_1,\ldots,Z_j$ be independent real variables, each supported in an
interval of length at most $b$, where $0<b<1$.
Put $W=\sum_i Z_i$ and $t=\sum_i\operatorname{Var}(Z_i)$.
Suppose $K>0$ satisfies
\[
                  1-\cos(2\pi x)\le Kx\quad(0\le x\le1).
\]
For every $w\in\R$ and every
$0<r\le2\sin^2(\pi b)/b^2$,
\begin{equation}\label{v9:harmonic:gain}
                     \E\{w-W\}\ge K^{-1}(1-e^{-rt}).
\end{equation}
In particular $K=8$ is valid without a finite scalar certificate.
The sharper value $K=23/5$ may be used with the scalar majorization in
Appendix~\ref{v9:count:appendix}.
\end{lemma}
\begin{proof}
Let $Z_i'$ be an independent copy. Since $|Z_i-Z_i'|\le b$ and sine is
concave on $[0,\pi]$, the chord from $0$ to $\pi b$ gives
\[
 \sin^2\bigl(\pi(Z_i-Z_i')\bigr)
 \ge\frac{\sin^2(\pi b)}{b^2}(Z_i-Z_i')^2.
\]
Consequently, writing $v_i=\operatorname{Var}(Z_i)$,
\[
 |\E e^{2\pi iZ_i}|^2
 =1-2\E\sin^2\bigl(\pi(Z_i-Z_i')\bigr)
 \le1-\frac{4\sin^2(\pi b)}{b^2}v_i
 \le e^{-2rv_i}.
\]
The intermediate right side is nonnegative because it bounds the
nonnegative left side. Taking square roots and multiplying, using
independence, gives $|\E e^{2\pi iW}|\le e^{-rt}$. Periodicity then gives
\[
 K\E\{w-W\}\ge1-\Re\E e^{2\pi i(w-W)}
                \ge1-|\E e^{2\pi iW}|\ge1-e^{-rt}.
\]
No centering or continuous distribution assumption was used. To verify
$K=8$, the function $y-1+\cos y$ is nondecreasing for $y\ge0$, since its
derivative is $1-\sin y\ge0$. Thus $1-\cos(2\pi x)\le2\pi x<8x$ for
$x>0$; $\pi<4$ also follows directly from
$\pi=4\int_0^1(1+y^2)^{-1}\,dy<4$. Equality at $x=0$ is harmless.
\end{proof}

The sine argument applies throughout $b<1$; the stronger restriction
$b<1/2$ is unnecessary for this step. For example,
$b=3/4$ permits the exact rate $r=16/9$.
This extends only the harmonic component. It does not extend other
denominator, binomial, support or noise gates. In the whole-count
notation a sufficient acquired span is $b=A+d_0$; it must be a bound on
each actual conditional residual, not on its variance.

For example, on a phase interval $[u_l,u_h]$, the individual
first-mark denominators are positive when
$1-2d_0-u_h>0$ and $1-A+u_l>0$, with $\zeta=1-d_0$.
On the full phase interval $[-1/2,1/2]$, these become $d_0<1/4$
and $A<1/2$. They imply $A+d_0<3/4<1$, so the extended harmonic
lemma adds no further span restriction there. The optional common-mark
bound has additional gates. None of these inequalities alone establishes
the final score, and a larger span gives a smaller harmonic rate that must
be used in that score.

\begin{proposition}[The boundary is a real obstruction]\label{v9:harmonic:boundary}
No strictly positive uniform rate in \eqref{v9:harmonic:gain} is possible for all
independent positive-variance summands of support span at most one.
Even a uniformly positive rate as $b\uparrow1$ is impossible.
\end{proposition}
\begin{proof}
For $0<\varepsilon<1$, let $Z$ equal $\varepsilon-1$ with probability
$\varepsilon$, and $\varepsilon$ otherwise. It is centered, has positive
variance $\varepsilon(1-\varepsilon)$ and span one.
For $j$ independent copies, $W=j\varepsilon-M$ with integer-valued
$M$, so $\{j\varepsilon-W\}=0$ identically. The right side of
\eqref{v9:harmonic:gain} is positive for every $r>0$ and finite $K>0$.

For the boundary limit, the centered variable taking values $\pm b/2$
with equal probabilities has variance $b^2/4$ and characteristic modulus
$|\cos(\pi b)|$. As $b\uparrow1$, this modulus tends to one while the
variance tends to $1/4$. A positive uniform exponential decay rate would
contradict these limits. Directly for \eqref{v9:harmonic:gain}, choose $w=-b/2$:
then $\E\{w-Z\}=(1-b)/2\to0$, also contradicting a positive uniform
rate with finite fixed $K$. In fact $1-|\cos(\pi b)|$ vanishes quadratically
in $1-b$, consistently with the rate in Lemma~\ref{v9:harmonic:span}.
\end{proof}

\subsection{Why small defect alone cannot replace confinement}
The obstruction can be realized by the actual conditional laws of a
centered, variance-one iid marginal with a large central gap. It is not
merely an optimizer of a moment relaxation.

For rational $\varepsilon,t>0$, set
\begin{equation}\label{v9:harmonic:family}
 q=\frac{1+\varepsilon-\varepsilon^2}{1+t^2},\quad
 p=1-q,\quad r=tq,
 \qquad 0<\varepsilon<p<1.
\end{equation}
Let $Y$ take the values $\varepsilon-1,\varepsilon,1$ with respective
probabilities $q\varepsilon,q(1-\varepsilon),p$, and put $X=(Y-p)/r$.
Direct calculation gives
\[
 \E Y=p,\quad
 \operatorname{Var}(Y)=q\{p+\varepsilon(1-\varepsilon)\}=r^2.
\]
Therefore $X$ is exactly centered and standardized. The first two atoms
are negative, and the third is positive. In the actual sign split,
$\ell=1/r$, $U_1=0$, and
\[
 U_0=\begin{cases}\varepsilon-1,&\text{probability }\varepsilon,\\
                  \varepsilon,&\text{probability }1-\varepsilon.
       \end{cases}
\]
Thus the actual conditional amplitude is $1-\varepsilon$, its span is
one, the normalized central-gap increment is $c_{\rm gap}=1-\varepsilon$,
and $\lambda=nq\varepsilon(1-\varepsilon)$.
In particular the central gap is wide even when the harmonic estimate
fails completely.

For a fixed interior count $k$, the actual conditional sum has the law
\[
 V_k=(n-k)\varepsilon-\operatorname{Bin}(n-k,\varepsilon).
\]
Thus $\E\{w-V_k\}=0$ when $w=(n-k)\varepsilon$, despite its positive
variance. At $n=42$, $k=17$, sufficiently small $\varepsilon$ makes
$u=\eta=w/2$ a legal marked phase, with $|u|\le1/2$ and $0<\eta\le1$.
Here $w=u+\eta$ is precisely the phase used by the marked-count argument.

\begin{corollary}[No moment-only amplitude implication]\label{v9:harmonic:arbitrary}
For every $\delta>0$ there is an actual centered variance-one finite law
supported in $[-96/25,96/25]$, with positive conditional noise,
$D<\delta$ and $\lambda<\delta$ at $n=42$, whose negative conditional
residual span is exactly one and whose fractional-part gain vanishes
at a legal marked phase. The normalized central-gap increment can be
arbitrarily close to one.
\end{corollary}
\begin{proof}
At $\varepsilon=0$ the positive-weight part of \eqref{v9:harmonic:family} is
the standardized two-point law at $-t,1/t$, with positive probability
$p=t^2/(1+t^2)$. Its defect is
\[
 D_B=\frac{(6+2\sqrt{10})(p-p_*)^2}{\sqrt{p(1-p)}},
 \qquad p_*=(4-\sqrt{10})/2.
\]
This follows by substituting $\rho_B,\gamma_B$ and $J_B=t+1/t$.
Choose positive rational $t$ arbitrarily near
$\sqrt{p_*/(1-p_*)}$. Then $D_B$ is arbitrarily small.
For this fixed $t$, let positive rational $\varepsilon$ tend to zero.
All moments are continuous. The extra outer-gap weight tends to zero,
so the definition of $J$ in\eqref{ms:physical-defect-definition} gives $D\to D_B$ as well;
$\lambda\to0$ while the conditional span remains one.
At the limiting value of $t$, the outer atom tends to
$-(1+2t^2)/t$. Here $0.84<t<0.85$, so its magnitude is less than
$2(0.85)+1/0.84<3$. The support is therefore within
$[-96/25,96/25]$ in a neighborhood. The sign, phase and positive-noise
conditions also hold for all sufficiently small $\varepsilon$.
\end{proof}

These laws are not asserted to satisfy a maximizing contact inequality,
nonnegative fixed-target excess, or all the signed-acquisition outputs.
They do not refute the sharp conjecture or the conditional local comparison.
They refute the proposed implication from small moment resources and an
actual central gap alone to a positive fractional-part gain or a small
conditional amplitude. Contact-based confinement or a quantitatively
proved replacement for harmonic nonresonance remains necessary for that proposed shortcut.

